% ============================================================
%  Harald Høffding, »Om Religionsfilosofiens Opgave og Fremgangsmaade«
%  Oversigt over det Kgl. Danske Videnskabernes Selskabs Forhandlinger
%  1900, Nr. 5, pp. 411--421.
%  Communicated at the meeting of 2 November 1900.
%  TRANSCRIPTION (Danish)
%
%  Source: scan.pdf in this directory (gitignored), 11 pages, from
%          publ.royalacademy.dk --
%          http://publ.royalacademy.dk/backend/web/uploads/2019-09-04/
%          AFL%203/O_1900_00_00_1900_4579/O_1900_18_00_1900_4573.pdf
%          400 ppi RGB JPX, 2306--2348 x 3650 px per page.
% ============================================================
%
%  PAGE MAP -- VERIFIED BY EYE ON ALL ELEVEN LEAVES (30 Aug 2026).
%    PDF page N == printed page N + 410, uniformly; pp. 411--421.
%    Printed p. 411 carries NO journal numeral: it is the opening page of the
%    article and bears the masthead instead. Its identity is fixed from p. 412
%    in the verso head above, and from the article's own citation.
%    Running heads: verso = numeral left, "Harald Høffding." centred;
%    recto = "Om Religionsfilosofiens Opgave og Fremgangsmaade." centred
%    (letterspaced), numeral right. Heads are not transcribed.
%    See pagemap.py (untracked) for the machine-readable map.
%
%  A SECOND FOLIO RUNS THROUGH THE ITEM. Every leaf carries, centred at the
%    foot beneath the text block, the offprint (Særtryk) pagination 1--11,
%    running alongside the journal's 411--421. On p. 421 it sits at the normal
%    foot line although the text stops halfway down the page, which is what
%    shows it to be a folio and not a catchword. It is furniture of the
%    separate issue, not of the text, and is NOT transcribed; it is recorded
%    here because a reader collating this edition against the scan will
%    otherwise meet two numbers on every page and not know which is which.
%
%  SIGNATURES. "27" at the foot of p. 411 and "27*" at the foot of p. 413 --
%    sheet signature and star signature of the parent volume, leaves 1 and 3
%    of the gathering, set right of centre and inset from the right margin.
%    Nowhere else in the article. ABBYY reads the star as a prime (27'); it is
%    an asterisk. Not transcribed. The foot of p. 411 also carries, at far
%    left, the volume line "D. K. D. Vid. Selsk. Overs. 1900."
%
%  TYPESETTING OF THE ORIGINAL.
%    Antiqua throughout; no Fraktur anywhere in the item.
%    NO EMPHASIS IN THE BODY AT ALL. Every word on all eleven pages was tested
%      for anomalous per-character advance width (Sperrsatz detection) and the
%      only hits are the small caps of the p. 411 masthead. There is no italic,
%      no letterspacing and no bold in the running text. Two passages read as
%      emphasised and are not -- p. 416 "religiøs Erfaring og religiøs Tro"
%      and p. 418 "det religiøse Axiom" -- where the apparent spacing is loose
%      justification. Consequently this file contains no \emph{} in the body,
%      and any that appears later is a defect.
%    Display type: the title and the author line are bold antiqua; the masthead
%      is letterspaced small caps.
%    The opening word "Ved" (p. 411) begins with an oversized initial V, about
%      two line-heights, baseline-aligned and raised rather than dropped.
%    NO FOOTNOTES anywhere in the item -- no *), no 1), no superscript figures,
%      no marginalia.
%    NO QUOTATION MARKS anywhere in the item. The one quoted saying, Heraclitus
%      on p. 420, is set unmarked after a colon. check.py therefore treats any
%      quotation mark in this file as suspect rather than balancing them.
%    Em-dashes are frequent and always spaced ( --- ) as printed.
%    Typographic apostrophes in genitives: "Augustin's Confessioner",
%      "Suso's og den hellige Teresa's Selvbiografier" (p. 416).
%    Acute accent: "paa én Gang" (p. 419).
%    The article closes on p. 421 with a decorative double rule. No date line,
%      no signature, no errata note.
%
%  CONVENTIONS. 1900 orthography exactly as printed -- "Fremgangsmaade",
%    "højeste", "Værdi", capitalised nouns, "aa" where printed. Not modernised.
%    Page markers "% --- p. N ---" sit at the line where the printed page
%    begins and are NOT surrounded by blank lines unless the printed text
%    really starts a new paragraph there. (Splicing a blank line in at a page
%    joint invents a paragraph: nine such false breaks had to be swept out of
%    pascal-kierkegaard, and check.py now flags the pattern here.)
%
%  PRINTER'S DEFECTS -- transcribed as printed, logged at the site, never
%    silently corrected. FOUR in the whole item, each verified on the image at
%    400 dpi, each carrying a "% sic:" comment at its place in the text below:
%      p. 415  "forholder sig med de religiøse Forestilliger" -- "Forestilliger"
%              for "Forestillinger", the n dropped.
%      p. 416  "at fastholde det højeste værdifulde, dets Oplevelser have ført
%              det til at kende." -- a dropped "som" before "dets".
%      p. 419  "Muligheden af et Princip, af hvilken paa én Gang" -- the neuter
%              "Princip" requires "hvilket".
%      p. 420  "Allerede dem gamle Heraklit har sagt" -- "dem" for "den".
%    ABBYY reproduced all four rather than normalising them, which is part of
%    why its layer was trusted as the primary witness here.
%
%    A METHOD WARNING, worth the space because it nearly cost a datum. The
%    p. 415 "Forestilliger" was filed by the recon pass as an ABBYY slip and
%    written into ocr.sh's sed table as a correction; the batch agent found the
%    page really prints it, and the sed line was silently repairing the
%    compositor. The line has been removed and a note left in ocr.sh. Every
%    entry in that table must be a slip checked against the image, because an
%    unverified "fix" there corrupts the edition where nothing will show it.
%
%  PAGE-JOINT HYPHENS. Two words are divided across a leaf and are closed up
%    here with "\-%", which keeps the printed division as a discretionary break:
%    pp. 416/417 "intel-"/"lektuel", pp. 417/418 "sammen-"/"lignende". The
%    suspended hyphen in "Følelses-" (p. 417) is a different thing and stands as
%    printed. No paragraph break falls at any page joint in this item: every one
%    of the ten joints lands inside a running paragraph, checked on the images.
%
%  QUOTE BALANCE: 0 marks of any sort in the whole item, by design of the
%    original. Not a running total to watch, unlike the Fraktur books.
% ============================================================
\documentclass[12pt,a4paper]{article}
\usepackage[utf8]{inputenc}
\usepackage[T1]{fontenc}
\usepackage{libertinus}
\usepackage{libertinust1math}
\usepackage[danish]{babel}
\usepackage[protrusion=true,expansion=true]{microtype}
\usepackage{setspace}
\usepackage[margin=1.2in]{geometry}
\usepackage{fancyhdr}
\usepackage{hyperref}

\hypersetup{
  pdftitle={Om Religionsfilosofiens Opgave og Fremgangsmaade},
  pdfauthor={Harald Høffding},
  hidelinks
}

\pagestyle{fancy}
\fancyhf{}
\fancyhead[LE,RO]{\thepage}
\fancyhead[RE]{\textit{Harald Høffding}}
\fancyhead[LO]{\textit{Om Religionsfilosofiens Opgave og Fremgangsmaade}}
\renewcommand{\headrulewidth}{0pt}

% The four section numbers stand alone on the page, centred, with white space
% above and below and no heading text of any kind. This is the whole of the
% article's internal division.
\newcommand{\snum}[1]{\bigskip\begin{center}#1\end{center}\nopagebreak\medskip}

\setstretch{1.15}

\begin{document}

% TITLE BLOCK AS PRINTED ON p. 411, top to bottom. The page marker for p. 411
% stands below this block rather than above it: the block is set by the
% skeleton, and the marker opens the transcribed running text of the same page.
%
%   masthead, two lines, letterspaced small caps;
%   a short centred rule;
%   the title, centred, bold antiqua;
%   "Af";
%   "Harald Høffding.", bold;
%   "(Meddelt i Mødet den 2. November 1900.)", smaller;
%   a short centred rule;
%   then "I." and the text, which opens with the oversized initial V of "Ved".
% The masthead is journal furniture rather than part of Høffding's text, but it
% is printed on the page and is set here, in small caps as printed, so that the
% opening leaf can be collated against the scan line for line.
\begin{center}
  {\footnotesize\scshape Oversigt over det Kgl.\ Danske Videnskabernes Selskabs\\
   Forhandlinger, 1900.\ Nr.\ 5.}

  \rule{4em}{0.4pt}

  \bigskip
  {\large\bfseries Om Religionsfilosofiens Opgave og Fremgangsmaade.}

  \medskip
  Af

  \smallskip
  {\bfseries Harald Høffding.}

  \smallskip
  {\small (Meddelt i Mødet den 2.\ November 1900.)}

  \medskip
  \rule{4em}{0.4pt}
\end{center}

\medskip

% --- p. 411 ---
% The opening word "Ved" is set with an oversized initial V, about two
% line-heights, baseline-aligned and raised rather than dropped. Display
% type, not emphasis; set plain here.
\snum{I.}

Ved Religionsfilosofi kan man forstaa to forskellige Bestræbelser. Man kan
mene en Filosoferen, en Tænkning, som udspringer af den religiøse Bevidsthed
som et Forsøg paa at vinde Klarhed over sig selv, om sin inderste Trangs
Retning og Genstand og om den Maade, paa hvilken denne Trang tilfredsstilles.
En saadan Eftertanke opstaar, naar Religionen ikke længere har Karakteren af
umiddelbar Følelse eller Instinkt, eller af tryg Hvilen i Traditionen, men den
bevæger sig dog indenfor selve Religionens Forudsætninger; Religionen er
Grundlag og Udgangspunkt, ikke egentligt Genstand. Jeg vil hellere kalde en
saadan Bestræbelse religiøs Tænkning end Religionsfilosofi. Nogle af de
største Aander, som have levet, have givet Bidrag til denne Art Tænkning.
Saaledes Augustinus, Middelalderens Mystikere, Pascal, Jacobi, Kierkegaard.
--- Ved Religionsfilosofi kan man ogsaa forstaa, og forstaar jeg her:
Filosoferen, Tænkning over Religionen som et givet psykologisk og historisk
Fænomen. Religionen er her Objekt, ikke Subjekt. Hin religiøse Tænkning falder
selv ind under det, som her er Objektet, idet den paa mange Maader stiller
Religionen i et klarere Lys, drager Konsekvenser
% --- p. 412 ---
og aabner Indblik, som den ureflekterede religiøse Bevidsthed ikke frembyder.

En Religionsfilosofi i den Betydning, i hvilken jeg tager Ordet, forudsætter,
at der existerer et religiøst Problem. Og dette forudsætter igen, at
Religionen ikke er den eneste Form for aandelig Kultur, og tillige, at flere
forskellige religiøse Standpunkter gøre sig gældende i større eller mindre
Modsætning til hverandre. Religionsfilosofien er komparativ: den søger ved
Sammenligning af Religionen med andre aandelige Bestræbelser saavel som ved
Sammenligning mellem forskellige religiøse Retninger indbyrdes at finde
Religionens egentlige Væsen og Princip, --- om muligt at formulere dette i en
enkelt Sætning, som de forskellige Religioner og religiøse Standpunkter i mere
eller mindre konsekvent og udpræget Form og i forskellige Iklædninger søge at
hævde. En saadan Sætning maa naturligvis have en vis abstrakt Form, og i denne
Form existerer den kun i Religionsfilosofien, ligesom f.\ Ex. Lovene for
Forestillingernes Association kun existere i Psykologien, medens det dog er
ifølge dem, at den praktiske Bevidstheds Forestillinger forbinde sig.

Et religiøst Problem kan ikke opstaa til enhver Tid. Det forudsætter en
Differentiation eller en Arbejdsdeling paa det aandelige Omraade. Det
existerer ikke i Religionens klassiske Tider, --- hverken i de store
Begyndelsestider, hvor den religiøse Interesse umiddelbart og fuldstændigt
lægger Beslag paa al Energi, eller i de store Organisationstider, hvor andre
aandelige Bestræbelser vel gøre sig gældende, men hvor de alle underordne sig
Religionen og fra denne hente deres sidste Kriterier og Regulativer. I
Religionens klassiske Tider kan der være religiøs Tænkning, men ingen egentlig
Religionsfilosofi. Det religiøse Problem forudsætter, at Enheden og Harmonien
paa det aandelige Omraade er brudt eller dog anfægtet. Videnskab og Kunst,
Moral og Samfundsliv udvikle og organisere sig efter deres egne Love, paa
deres selvstændige Grundlag, og staa som selvstændige Aandsmagter overfor
Religionen. Og indenfor Religionen
% --- p. 413 ---
gøre sig tillige forskellige Retninger gældende, ja Forskellighederne angaa
maaske endog hvad der fra den ene eller den anden Side anses som det
væsentlige i Religionen.

Spørgsmaalet bliver da, hvilken Betydning Religionen under disse Forhold kan
have. Hvis det er Religionens Væsen at være den eneste Form for aandelig
Kultur, og hvis det ligger i Religionens Væsen at fremtræde under en eneste
bestemt Form, saa er det klart, at dens Tid er forbi, naar aandelig
Arbejdsdeling og Differentiation er fremtraadt. I Virkeligheden høres ogsaa en
saadan Paastand ofte, --- medens det paa den anden Side forsikres, at hele den
intellektuelle, æstetiske, moralske og sociale Udvikling, der er gaaet for
sig, aldeles ikke har ændret Religionens Stilling indenfor Aandslivet. At saa
modsatte Paastande kunne opstilles, godtgør, at her virkeligt er et Problem.

Det er ikke rimeligt, at dette Problem vil kunne løses. Vi kunne paa det
aandelige Omraade langt ufuldkomnere end paa det fysiske Omraade tydeligt
skelne de enkelte Elementer og Kræfter og finde Lovene for deres Samvirken,
saa at man klart skulde kunne afgøre, om Livsbetingelserne for en fortsat
Bestaaen af en aandelig Retning ville vedblive at være til Stede eller ikke.

Til Undersøgelse af dette Spørgsmaal vil det ikke nytte at stille Filosofien
som et færdigt System overfor Religionen. Det er Filosofien som Organ, som
alsidig Drøftelse af Aandslivet og dets Betingelser, der her er Brug for. Det
vil naturligvis kunne have sin Interesse at undersøge, hvorledes en fuldt
udviklet filosofisk Verdensanskuelse forholder sig til Religion. Men for det
første har et saadant System mere Karakteren af Kunst end af Videnskab; det
forholder sig til den undersøgende Filosofi som Historieskrivningen forholder
sig til historisk Kritik. Og for det andet vil et filosofisk System kun have
Værdi, naar det allerede under sin Tilblivelse og Forberedelse har taget
tilbørligt Hensyn til et saadant aandeligt Faktum som Religionen er.
% --- p. 414 ---
Det vil da være naturligt ved Drøftelsen af det religiøse Problem at anvende
Filosofien i dens undersøgende og forberedende Skikkelse. Filosofi er i denne
Forstand en kritisk og sammenlignende Disciplin, der opsøger Bevidsthedslivets
Grundformer og Grundforudsætninger og Vilkaarene for deres Bestaaen og deres
Gyldighed. Dens Metode er væsentligt analytisk-genetisk, idet den fra de
færdige Værker og Former gaar tilbage til de Kræfter og Elementer, af hvilke
de ere fremgaaede. En Undersøgelse af Religionen af denne Art vil
hensigtsmæssigt kunne anlægges saaledes, at den bestaar af en
erkendelsesteoretisk, en psykologisk og en etisk Del. --- Jeg skal forsøge at
give en sammentrængt Oversigt over en saadan Undersøgelse.

\snum{II.}

I Religionens klassiske Tider, i hvilke den var et og alt, afgav de religiøse
Forestillinger tilfredsstillende Forklaring og Forstaaelse af Tilværelsen. Der
har nu udviklet sig en Trang til Forstaaelse, som ikke fyldestgøres ved disse
Forestillinger. Det enkelte Fænomen, den enkelte Begivenhed i den aandelige og
i den materielle Verden forstaas (efter dette nye Begreb om Forstaaelse) ved
at bringes i den nøjeste, kontinuerlige Sammenhæng med andre Fænomener og
Begivenheder, der hver for sig ligesaa vel kunne paavises i Erfaringen som de,
hvis Forklaring man spurgte om. Det fordres, at Aarsagen ligesaa vel som
Virkningen skal være en Erfaring. Mere og mere anerkendes det nu (skønt det
har kostet haarde og gentagne Kampe), at det ikke er Religionens Sag at yde en
Forstaaelse af denne Art. Og naar de to Forklaringsmaader støde sammen ved et
og samme Fænomen, bliver Spørgsmaalet, hvorledes de skulle forenes, og om de i
det hele kunne det. Saa meget synes i det hele dog at være almindeligt
indrømmet, at de religiøse Forestillinger ikke have deres Betydning og deres
Værdi derved, at de skulle yde fuld Tilfredsstillelse af Erkendelsestrangen.
Det
% --- p. 415 ---
er da sandsynligt, at de religiøse Forestillinger skylde deres Opstaaen og
Bestaaen andre Sider af Bevidsthedslivet end den intellektuelle.

Man kunde mene, at det kun var de enkelte specielle Fænomener, Videnskaben
kunde forklare, medens de afsluttende Spørgsmaal, de sidste Gaader maatte
søges besvarede i Religionen. Men hvis det nu viser sig ved en nærmere kritisk
Undersøgelse, at Religionen ligesaa lidt som Videnskaben formaar at løse de
sidste Gaader, give os virkeligt afsluttende og konsekvente Begreber, saa
vender Spørgsmaalet om de religiøse Forestillingers Betydning tilbage med
fornyet Kraft. Det viser sig (som det her vilde være for vidtløftigt at
godtgøre i det enkelte), at religiøse Forestillinger væsentligt skyldes
Analogi, uden at de Analogier, der her ligge til Grund, kunne formes og
anvendes paa en saa klar og gennemført Maade som ved de Analogier, den
videnskabelige Tænkning paa forskellige Omraader gør Brug af. Dette er ingen
Tilfældighed; ti som det
% sic: printed "Forestilliger" (no n), verified at 400 dpi. The sense
% requires "Forestillinger". NOT an ABBYY slip: ocr.sh's sed table silently
% "corrects" this form, which is why the recon pass took it for one.
forholder sig med de religiøse Forestilliger, forholder det sig ogsaa med de
Tanker, med hvilke den spekulative Filosofi i de forskellige metafysiske
Systemer har forsøgt at give en Afslutning paa Forstaaelsen af Tilværelsen.
Der kommer her et Punkt, hvor Tanken afløses af Poesien. De religiøse
Forestillinger have jo ogsaa gennemgaaende en billedlig Karakter. Det er store
Symboler, i hvilke aandelige Tilstande af anden Art og Omfang end de rent
intellektuelle, have søgt og fundet Udtryk. --- Den erkendelses-teoretiske
Religionsfilosofi fører paa dette Punkt naturligt over i den psykologiske.

\snum{III.}

Denne Del af Religionsvidenskaben vil mere og mere blive den vigtigste, og det
er den, der for mig har frembudt den største Interesse. Vi føres her bort fra
Striden mellem Dogmatik og Kritik tilbage til Sjælelivet i dets inderligste og
mest
% --- p. 416 ---
% sic: printed "det højeste værdifulde, dets Oplevelser have ført det til at
% kende." -- a dropped "som" before "dets". Verified at 400 dpi.
koncentrerede Former. Religion kan ikke laves eller konstrueres. Den voxer ud
af Livet, udspringer af Menneskets Grundstemning under Livets Kamp, af dets
Villie til at fastholde det højeste værdifulde, dets Oplevelser have ført det
til at kende. Vi lære den ikke altid saa godt at kende fra denne Side ved
Religionshistoriens Hjælp, saa betydningsfuld denne ellers naturligvis er. I
Religionshistorien faar man ikke saa let de sjælelige Kræfter i deres
oprindelige Virken at se. Den beskæftiger sig mere med Religionens Værker end
med Religionen selv, --- mere med de ydre Anledninger og Betingelser end med
det indre Liv, --- mere med de store Typer, de for store Menneskegrupper
fælles Former end med de individuelle Processer. Studiet af den individuelle
Religiøsitet, saaledes som denne kan kendes af religiøse Personligheders
Biografier, særligt Selvbiografier, er derfor af største Betydning for
Religionspsykologien. Augustin's Confessioner, Suso's og den hellige Teresa's
Selvbiografier ere psykologiske Kildeskrifter af første Rang. Naturligvis er
der, og det særligt paa det religiøse Omraade, en nøje Sammenhæng mellem
Individ og Samfund, og derfor er en Undersøgelse af Forholdet mellem den
individuelle Religiøsitet og den religiøse Tradition et af de vigtigste
Kapitler i Religionspsykologien. Den individuelle Religiøsitet hører desuden
ogsaa til Religionshistorien, skønt denne hidtil saa godt som udelukkende har
beskæftiget sig med de store Typer og set bort fra den mangfoldige Maade, paa
hvilken de nuanceres hos de enkelte Individer. Og dog er det egentligt først
gennem disse Nuancer, at de virkende Kræfter og Elementer træde tydeligt frem.

Der er to Begreber, det for Religionspsykologien især drejer sig om at
bestemme: religiøs Erfaring og religiøs Tro.

Den religiøse Erfarings Ejendommelighed finder jeg i, at medens andre
Erfaringer enten angaar den givne Virkelighed, hvilket gælder de Erfaringer,
vi gøre ved de forskellige Sanser, eller den Værdi, denne Virkelighed har for
os, hvilket gælder de Oplevelser, der skyldes organisk Livsfølelse, æstetisk,
intel\-%
% PAGE-JOINT HYPHEN. p. 416 ends "intel-" and p. 417 opens "lektuel", inside a
% running paragraph. Set as "intel\-%" so the two halves close up to
% "intellektuel" with the printed division kept as a discretionary break.
% --- p. 417 ---
lektuel og etisk Følelse, saa angaar den religiøse Erfaring eller den religiøse
Følelse Forholdet mellem Værdi og Virkelighed. Den er bestemt ved det
værdifuldes Skæbne i Virkelighedens Verden. --- Den religiøse Tros
Ejendommelighed finder jeg i, at den under de virkelige Begivenheders Gang
fastholder Værdiens Bestaaen. Den er en dyb Tillid, en Kontinuitet i Sindet
under alle Svingninger. Dens objektive Indhold er ogsaa en Kontinuitet: dens
Grundantagelse er --- antropomorfistisk udtrykt --- en Trofasthed i
Tilværelsen, som lægger sig for Dagen i, at Værdierne bestaa, saa meget end
deres specielle Former kunne skifte.

Denne psykologiske Karakteristiks Rigtighed søges godtgjort ad forskellige
Veje. For det første er det muligt ud fra de angivne Karaktermærker at give en
Beskrivelse af en Følelses- og Villiesretning, der naturligt maa udvikle sig
% The hyphen in "Følelses- og Villiesretning" is a suspended hyphen printed at
% a line end; it is retained, not a word-break hyphen.
under menneskelige Livsvilkaar og i Kraft af de almindelige psykologiske Love.
Det er en psykologisk mulig Tilstand. For det andet vil man ved at benytte
denne Beskrivelse og kalde det, den beskriver, Religion, paa en klar og tydelig
Maade kunne paavise den religiøse Følelses Forskel fra og Forhold til andre
Arter af Følelse, med hvilke den er beslægtet. Og for det tredie viser en
Undersøgelse af de vigtigste historisk givne Typer for Religion, at de nævnte
Karaktermærker ere de vigtigste og de gennemgaaende.

Sætningen om Værdiens Bestaaen kan siges at være det religiøse Axiom. Den
udtrykker en Grundtanke, som den religiøse Bevidsthed fastholder under højst
forskellige Former, og som i Debatten mellem forskellige religiøse Standpunkter
mere eller mindre tydeligt fremtræder som Ideal og Maalestok. Forskellen mellem
religiøse Standpunkter beror dels paa de Værdier, hvis Bestaaen der tros paa,
dels paa den Klarhed, hvormed Tanken om Værdiernes Bestaaen fremtræder og
fastholdes, dels paa de forskellige Maader, paa hvilke den udtrykkes. For den
enkelte bestemte Religion er vel den bestemte Form, i hvilken den udtrykker
Axiomet, det væsentlige. Men for den sammenlignende
% The word "sammenlignende" is broken across the page joint: p. 417 ends
% "sammen-" and p. 418 opens "lignende". Rejoined here, so the whole word
% stands before the p. 418 marker.
% --- p. 418 ---
Undersøgelse gælder det at finde det typiske, Temaet, der i det enkelte kan
varieres paa mange Maader.

Ved Hjælp af den saaledes begrundede Hypotese bliver det nu muligt at drøfte
det Spørgsmaal, som er Kærnen i det religiøse Problem, nemlig om Religionens
Bestaaen under de ændrede Vilkaar, Arbejdsdelingen paa det aandelige Omraade
medfører. Hvis Hypotesen er rigtig, vil der være at skelne mellem Religionens
Form og dens Kærne. Formerne --- Myter og Dogmer --- kunne skifte og Kærnen dog
bestaa. Derfor er den rent intellektuelle Kritik tilsidst af saa underordnet
Betydning. Fordi visse religiøse Forestillinger falde bort, falder ikke det
væsentlige Motiv til disse Forestillingers Dannelse bort. Det kan ytre sig i
nye Virkninger, ligesom Trangen til at finde Aarsager ikke falder bort, fordi
visse Aarsagsforklaringer have vist sig at være ugyldige. Religionens Tid vilde
kun være forbi, dersom det kunde bevises, at det var umuligt at fastholde
Antagelsen af Værdiens Bestaaen.

Det religiøse Axiom kan ikke udledes af Videnskaben. Den videnskabelige
Tænkning har fuldført sit Arbejde, naar den har paavist saa mange Forhold af
Identitet, Rationalitet og Kausalitet, som det til given Tid er den muligt. Det
Verdensbillede, som derved fremkommer, beholder sin Gyldighed, hvorledes saa
end Værdiernes Skæbne i Virkelighedens Verden bliver. Men det religiøse Axiom
behøver paa den anden Side ikke at føre til Strid med det videnskabelige
Verdensbillede, skønt dette faktisk ofte har været Tilfældet. Den Mulighed er
ikke udelukket, at det kunde udtrykkes paa en Maade, der udelukkede en saadan
Strid. Eller selve det virkelige Verdensbilledes Former og Love kunde netop
staa som de Veje og Midler, der gøre Værdiens Bestaaen mulig. Det vilde være
Idealet, om Værdisynspunktet og Virkelighedssynspunktet tilsidst kunde falde
sammen. Det Spørgsmaal kan fra den kritiske Filosofis Standpunkt ikke holdes
tilbage, om ikke den Forskel, vi gøre og maa gøre mellem Værdi og Virkelighed,
blot er subjektiv og menneskelig. Denne
% --- p. 419 ---
Forskel kunde jo bero paa, at baade vort Begreb om Værdi og vort Begreb om
Virkelighed er ufuldkomment, uafsluttet og uafslutteligt. Vi kunne hverken
udlede Værdien af Virkeligheden eller Virkeligheden af Værdien. Vor Tanke
henledes da til Muligheden af et Princip, af hvilken paa én Gang Værdiens
% sic: printed "af et Princip, af hvilken" -- the neuter "Princip" requires
% "hvilket". Reading verified at 400 dpi; the glyphs are unambiguous.
Bestaaen og Virkelighedens faste Sammenhæng kunde udspringe. Men det er en
Tanke, der ikke kan videnskabeligt gennemføres.

Der vil stadigt være Stof nok til Strid paa det religiøse Omraade. Der vil
kunne være Strid om, hvilke Værdier det er, hvis Bestaaen der skal tros paa, og
der vil være Strid om de rette Former for Troen paa disse Værdiers Bestaaen.
Men det fundamentale religiøse Stridspunkt vil være det, om selve det religiøse
Axiom kan holdes oppe. Mere end en Mulighed heraf kan teoretisk ikke godtgøres,
dersom det er rigtigt, at al Tro er Livets Sag. Men Livet er maaske rigere og
frugtbarere end bornerede Tvivl og bornerede Dogmer antage.

\snum{IV.}

Den etiske Religionsfilosofi danner naturligt Afslutningen af den hele
Undersøgelse. Den undersøger to Hovedspørgsmaal, nemlig hvorvidt Religion er
Grundlag for Etik, og hvilken Betydning Religionen --- selv om dette ikke
skulde være Tilfældet --- har som en Form af aandelig Kultur. Begge Spørgsmaal
kunne vinde i Klarhed ved Hjælp af den Hypotese, at Tro paa Værdiens Bestaaen
er det væsentlige i Religionen.

De Værdier, hvis Bestaaen Mennesket tror paa, maa det først selv kende af
Erfaring. Naar etiske Ideer fremtræde i religiøs Form, f. Ex. naar Guderne
fremtræde med etiske Egenskaber, saa er Forudsætningen, at disse Ideer eller
Egenskaber i selve det menneskelige Liv have godtgjort deres Værdi. Kun saa kan
Mennesket forbinde nogen Mening med dem, naar de benyttes til at udtrykke det
guddommelige. Religionshistorikerne lægge derfor stor Vægt paa Udviklingen af
Menneskenes etiske
% --- p. 420 ---
Forestillinger som Betingelse for den vigtigste Modsætning indenfor
Religionshistorien, Modsætningen mellem Naturreligion og etisk Religion. Baade
historisk set og rent logisk set har det etiske derfor Prioritet i Forhold til
det religiøse. Og dette træder særligt frem derved, at religiøse Debatter
tilsidst føre tilbage til etiske Modsætninger, til en forskellig Vurdering, ---
altsaa til en Forskel i Henseende til de Værdier, hvis Bestaaen der tros paa.

Men selv om dette er rigtigt, er dermed Forholdet mellem Religion og Etik ikke
udtømt. Der maa endnu fra et rent etisk Standpunkt spørges: hvilken Værdi har
det at tro paa Værdiens Bestaaen? Det er Spørgsmaalet om Religionens Betydning
som Form for aandelig Kultur. Behandlingen af dette Spørgsmaal kræver en
indgaaende psykologisk og sociologisk Undersøgelse. Psykologisk set har en
saadan Tro Betydning ved at styrke Villien og ved den Udvidelse af Horizonten,
den fører til. Troen paa Værdiens Bestaaen lader os, selv om vi ikke i det
enkelte kunne forfølge Værdimetamorfoserne, anskue Tilværelsen som et stort
System af potentielle Værdier, der for en Del gennem menneskeligt Arbejde
skulle omsættes til aktuelle Værdier. Gennem sin inderste Rod hænger her den
etiske Bevidsthed sammen med den religiøse. Troen paa Værdiens Bestaaen faar i
Etiken en lignende Betydning som Ideer, Antecipationer og Hypoteser have paa
det teoretiske Omraade. En ængstelig Klamren til givne Erfaringer udelukker let
fra at gøre nye Erfaringer. Allerede dem gamle Heraklit har sagt: Naar du ikke
% sic: printed "Allerede dem gamle Heraklit" -- "dem" for "den". Verified at
% 400 dpi. The Heraclitus saying that follows is set unmarked after the colon;
% no quotation marks in the printed text.
haaber, vil du ikke udfinde det, der ikke kunde haabes. Sociologisk har
Religionen sin store Betydning gennem sin Organisation og gennem den
filantropiske Virksomhed, den fremkalder.

Et mere eller mindre skarpt Modsætningsforhold mellem Religion og Etik
fremkommer ikke blot, naar Religionen opstiller Bud eller Forbud, som etisk set
ikke kunne anerkendes. Den mest karakteristiske Modsætning fremkommer, naar
Troen paa Værdiens Bestaaen selv bliver en Hindring for Arbejdet paa
% --- p. 421 ---
at opdage og frembringe Værdier ved at lægge Beslag paa Kraft og Tid, der
ellers kunde være brugt i dette Arbejdes Tjeneste. Dersom der er dem, hvis hele
Kraft gaar op i et saadant Arbejde, bliver det Spørgsmaalet, om det forringer
deres Personligheds Værdi, at de ikke --- i hvert Tilfælde ikke med fuld
Bevidsthed --- tro paa Værdiens Bestaaen i nogensomhelst Form. Dette er maaske
den mest afgørende Form, i hvilken det religiøse Problem kan træde frem.

Det synes da, som om den religionsfilosofiske Opfattelse, der her er antydet i
de store Træk, ikke blot har psykologisk og historisk Begrundelse, men ogsaa
formaar at afgive vigtige Synspunkter for Kampen paa det religiøse Omraade.
% End of the article. The text stops here, about a third of the way down
% p. 421; a short centred decorative double rule follows, then the offprint
% folio "11" at the foot. The rule is not set. There is no date line and no
% signature.

\end{document}
